% Part: computability
% Chapter: computability-theory
% Section: ce-closed-cup-cap

\documentclass[../../../include/open-logic-section]{subfiles}

\begin{document}

\olfileid{cmp}{thy}{clo}

\olsection[اجتماع و اشتراک مجموعه‌های م.ب.]{بستگی مجموعه‌های
  محاسبه‌پذیرِ برشمردنی تحت اجتماع و اشتراک}

قضیهٔ زیر چند ویژگی بستاری مجموعه‌های محاسبه‌پذیرِ برشمردنی را به
دست می‌دهد.

\begin{thm}
فرض کنید $A$ و~$B$ محاسبه‌پذیرِ برشمردنی باشند. آنگاه $A\cap B$ و
$A\cup B$ نیز چنین‌اند.
\end{thm}

\begin{proof}
\olref[eqc]{thm:ce-equiv} اجازه می‌دهد از توصیف‌های گوناگون مجموعه‌های
محاسبه‌پذیرِ برشمردنی استفاده کنیم. برای توضیح مطلب چند برهان متفاوت
می‌آوریم.

برای برهان نخست فرض کنید $A$ با تابع محاسبه‌پذیر~$f$ و $B$ با تابع
محاسبه‌پذیر~$g$ برشمرده شود. بگذارید
\begin{align*}
h(x) & = \umin{y}{(f(y) = x \lor g(y) = x)} \text{ و}\\
j(x) & = \umin{y}{(f((y)_0) = x \land g((y)_1) = x)}.
\end{align*}
آنگاه $A\cup B$ دامنهٔ $h$ و $A\cap B$ دامنهٔ~$j$ است.

\begin{explain}
از نظر محاسباتی چنین رخ می‌دهد: با داشتن روش‌هایی که $A$ و~$B$ را
برمی‌شمارند، می‌توانیم با جست‌وجوی~$x$ در یکی از دو برشماری، عضویت
$x$ در $A\cup B$ را نیمه‌تصمیم بگیریم؛ و با جست‌وجوی هم‌زمان~$x$ در
هر دو برشماری، عضویت آن در $A\cap B$ را نیمه‌تصمیم بگیریم.
\end{explain}

برای برهان دوم، باز فرض کنید $A$ با~$f$ و $B$ با~$g$ برشمرده شود.
بگذارید
\[
k(x) = \begin{cases}
f(x/2) & \text{اگر $x$ زوج باشد} \\
g((x-1)/2) & \text{اگر $x$ فرد باشد.}
\end{cases}
\]
آنگاه $k$ مجموعهٔ $A\cup B$ را برمی‌شمارد؛ اندیشه این است که $k$
به‌تناوب از برشماری‌های $f$ و~$g$ می‌گیرد. برشمردن $A\cap B$ دشوارتر
است. اگر $A\cap B$ تهی باشد، بدیهی است که محاسبه‌پذیرِ برشمردنی است.
وگرنه بگذارید $c$ عنصری دلخواه از $A\cap B$ باشد و $l$ را چنین تعریف
کنید:
\[
l(x) = \begin{cases}
f((x)_0) & \text{اگر $f((x)_0) = g((x)_1)$} \\
c & \text{در غیر این صورت.}
\end{cases}
\]
از نظر محاسباتی، $l$ زوج‌هایی از عناصر برشماری‌های $f$ و~$g$ را
مرور می‌کند و هر تطابقی را که بیابد بیرون می‌دهد؛ در غیر این صورت با
بیرون‌دادن~$c$ صرفاً درجا می‌زند.

برای برهان آخر، فرض کنید $A$ \emph{دامنهٔ} تابع جزئی~$m(x)$ و $B$
دامنهٔ تابع جزئی~$n(x)$ باشد. آنگاه $A\cap B$ دامنهٔ تابع جزئی
$m(x)+n(x)$ است.

\begin{explain}
از نظر محاسباتی، اگر $A$ مجموعهٔ مقادیری باشد که $m$ در آن‌ها متوقف
می‌شود و $B$ مجموعهٔ مقادیری که $n$ در آن‌ها متوقف می‌شود، آنگاه
$A\cap B$ مجموعهٔ مقادیری است که هر دو روش در آن‌ها متوقف می‌شوند.
\end{explain}

بیان $A\cup B$ به‌صورت مجموعه‌ای از مقادیر توقف دشوارتر است، زیرا
باید $m$ و~$n$ را به‌طور موازی شبیه‌سازی کرد. بگذارید $d$ نمایه‌ای
برای $m$ و $e$ نمایه‌ای برای $n$ باشد؛ یعنی $m=\cfind{d}$ و
$n=\cfind{e}$. آنگاه $A\cup B$ دامنهٔ تابع زیر است:
\[
p(x) = \umin{y}{(T(d,x,y) \lor T(e,x,y))}.
\]
\begin{explain}
از نظر محاسباتی، $p$ روی ورودی $x$ در پی محاسبهٔ متوقف‌شونده‌ای برای
$m$ یا $n$ می‌گردد و با یافتن هر یک متوقف می‌شود.
\end{explain}
\end{proof}

\end{document}
